\appendix
\section{Building the Book and Checking Its Examples}
\label{sec:book-build}

From the repository root, build the PDF and check the examples:

\begin{lstlisting}[language=bash]
make -C tutorial
make -C tutorial check
\end{lstlisting}

The PDF build runs \code{lualatex} three times to stabilize the contents,
page numbers, and citations.
It keeps build intermediates in \filepath{build/tutorial-build/} and
copies the finished document to \filepath{doc/tutorial.pdf}.
The build and checks create their scratch directories automatically;
no pre-existing contents of \filepath{build/} are required.
It requires SBCL for source-excerpt preparation and a LuaLaTeX installation
with the packages named in the preamble,
Libertinus Serif and Libertinus Sans for text, and Noto Sans Mono for code.
The native runner \filepath{tutorial/tests/run-tests.lisp} needs SBCL and
the current ABCL/R2 executable, built with \code{make} at the root;
no Python or external Lisp packages are needed.
The checks run example sessions in fresh processes and keep source copies,
translations and transcripts under \filepath{build/}.
Session inputs are in \filepath{tutorial/tests/sessions/}; expected
results are specified in \filepath{tutorial/tests/cases.lisp}.
Run \code{make -C tutorial reading-check} for the Lisp-reading examples,
\code{reflection-check} for reflective protocols, and
\code{group-check} for ownership and identity checks.

Interactive loading creates a neighboring \code{.lsp} file; use
\code{:use-source t} while editing to avoid cached translations.
Printed object identifiers are session-dependent, not stable expected outputs.

\section{Original Source Reading Map}
\label{sec:source-map}

This map covers all 34 original example files and nine library files.
Each entry points to the explanation of its role in a complete example,
not necessarily to a separate runnable program. Start with the linked
application or loading section before reading its supporting generators.
Consult the walkthrough's stated checks and limits: a source-map entry
is not a certification that every execution path works on this image.

\newcommand{\sourcefile}[3]{\filepath{#1} & Section~\ref{#2}: #3\\}

\paragraph{Individual examples and reflective mechanisms.}
\begin{longtable}{L{0.41\textwidth}L{0.49\textwidth}}
\textbf{Original file} & \textbf{Reading route and role}\\\hline
\endhead
\sourcefile{examples/misc/future-ex.abclr2}{sec:original-future}{Tree traversal, repeated replies, and the terminal marker.}
\sourcefile{examples/misc/inherit.abclr2}{sec:original-inherit}{Bird/emu application, inherited scripts, and reflective context.}
\sourcefile{examples/misc/built-ins.abclr2}{sec:builtin-reading}{Schematic environment, script, and queue protocols; not a runtime to load.}
\sourcefile{examples/misc/user-script.abclr2}{sec:user-script-reading}{Sender-sensitive script matching, dynamic installation, and separate application/evaluation replies.}
\sourcefile{examples/misc/MCP.abclr2}{sec:mcp-reading}{Meta-circular group manager, application metaobjects, and evaluator cases.}
\sourcefile{examples/misc/meta-meta.abclr2}{sec:meta-meta-reading}{Three levels of identities and the twice computation.}
\sourcefile{examples/misc/reify-args.abclr2}{sec:original-if}{Eager/lazy branch selection; Section~\ref{sec:interrupt} follows interruption and the separate control-flow experiment.}
\sourcefile{examples/misc/fault-tolerant.abclr2}{sec:retry}{Acceptance state, initialization, rejection, and retry with preserved reply/sender.}
\sourcefile{examples/misc/break.abclr2}{sec:break}{Breakpoint continuation, lexical environment navigation, and resume.}
\sourcefile{examples/misc/express-ex.abclr2}{sec:original-express}{Newer clock protocol, suspension context, atomic regions, and express-message routing.}
\sourcefile{examples/misc/express-old-ex.abclr2}{sec:original-express}{Legacy clock and reader/writer lock participation; separate fresh session.}
\end{longtable}

\paragraph{Priority control: read the application, then the execution policy.}
\begin{longtable}{L{0.41\textwidth}L{0.49\textwidth}}
\textbf{Original file} & \textbf{Reading route and role}\\\hline
\endhead
\sourcefile{examples/pcontrol/README.abclr2}{sec:pcontrol-load}{Historical entry instructions, argument shapes, and small inputs.}
\sourcefile{examples/pcontrol/start.abclr2}{sec:pcontrol-load}{Compile applications, then load the four machinery files in dependency order.}
\sourcefile{examples/pcontrol/fib.abclr2}{sec:pcontrol-fib}{Child creation, saved reply destination, and two-answer join.}
\sourcefile{examples/pcontrol/qsort.abclr2}{sec:pcontrol-qsort}{Partition self-messages, shared array, and completion join.}
\sourcefile{examples/pcontrol/pcontrol-group.abclr2}{sec:pcontrol-policy}{Connect priority metaobjects and four evaluator objects into one group.}
\sourcefile{examples/pcontrol/priori-meta.abclr2}{sec:pcontrol-policy}{Attach object priority to evaluation and announce dispatch activity.}
\sourcefile{examples/pcontrol/pqueue-eval.abclr2}{sec:pcontrol-policy}{Priority/FIFO work queues, priority associations, and cyclic evaluator routing.}
\sourcefile{examples/pcontrol/para-monitor.abclr2}{sec:pcontrol-monitor}{Logical clock, active/object counters, and output streams.}
\end{longtable}

\paragraph{Time Warp: separate application, histories, and evaluator scheduling.}
\begin{longtable}{L{0.41\textwidth}L{0.49\textwidth}}
\textbf{Original file} & \textbf{Reading route and role}\\\hline
\endhead
\sourcefile{examples/time-warp/README.abclr2}{sec:tw-loading}{Historical run sequence and virtual-time application interface.}
\sourcefile{examples/time-warp/tw-load.abclr2}{sec:tw-loading}{Ten-file compile/load sequence; distinguishes loaded machinery from alternatives.}
\sourcefile{examples/time-warp/tw-wash1.abclr2}{sec:tw-application}{Timed car/washer pairing, availability, finish sentinel, and controller join.}
\sourcefile{examples/time-warp/tw-wash0.abclr2}{sec:tw-loading}{Earlier ordinary-message application, without Time Warp or timed service completion.}
\sourcefile{examples/time-warp/tw-meta.abclr2}{sec:tw-recovery}{Ordinary/timed dispatch, stragglers, state restoration, and anti-messages.}
\sourcefile{examples/time-warp/tw-undo-env.abclr2}{sec:tw-histories}{Saved local binding frames and strict-before-cutoff restoration.}
\sourcefile{examples/time-warp/tw-undoable-pqueue.abclr2}{sec:tw-histories}{Pending/dequeued inputs, replay, cancellation, and active versus conditional alternatives.}
\sourcefile{examples/time-warp/tw-output-history.abclr2}{sec:tw-histories}{Record sent events and extract speculative output by send-time cutoff.}
\sourcefile{examples/time-warp/tw-group.abclr2}{sec:tw-evaluator}{Connect components, transfer scheduler queues, and select evaluator variants.}
\sourcefile{examples/time-warp/tw-eval.abclr2}{sec:tw-evaluator}{Expression interpretation and time/output context; Section~\ref{sec:tw-variants} compares NR/NS.}
\sourcefile{examples/time-warp/tw-meta-eval.abclr2}{sec:tw-evaluator}{Queue evaluator requests by LVT and count evaluation arrivals.}
\sourcefile{examples/time-warp/tw-pqueue.abclr2}{sec:tw-evaluator}{Priority and FIFO work queues; not the undoable application input queue.}
\sourcefile{examples/time-warp/tw-monitor.abclr2}{sec:tw-evaluator}{Rollback announcements and their counter.}
\sourcefile{examples/time-warp/tw-test.abclr2}{sec:tw-evaluator}{Independent countdown objects; the starter is not a completion join.}
\sourcefile{examples/time-warp/tw-scheduler.abclr2}{sec:tw-evaluator}{Unfinished fragment, not loaded by the application loader.}
\end{longtable}

\paragraph{Libraries: loading is not the same as installing a policy.}
\begin{longtable}{L{0.41\textwidth}L{0.49\textwidth}}
\textbf{Original file} & \textbf{Reading route and role}\\\hline
\endhead
\sourcefile{lib/future.abclr2}{sec:original-future}{Buffered reply destination, consumption, and explicit application completion; see also Section~\ref{sec:futures}.}
\sourcefile{lib/multicast.abclr2}{sec:multicast}{Evaluator installation, target templates, forwarding, and reply collection.}
\sourcefile{lib/heavy-meta.abclr2}{sec:guards}{Guard-aware script selection and candidate fallback.}
\sourcefile{lib/heavy-meta-old.abclr2}{sec:guards}{Historical reference only: use the regular variant with the current translator.}
\sourcefile{lib/express.abclr2}{sec:original-express}{Newer sender/receiver machinery and execution modes.}
\sourcefile{lib/express-old.abclr2}{sec:original-express}{Different legacy lock/count protocol; not an alias for the newer library.}
\sourcefile{lib/copy.abclr2}{sec:copy}{New identities versus shared lexical/script/reply references and state-copy restrictions.}
\sourcefile{lib/watch.abclr2}{sec:watch}{Outgoing-send evaluator wrapper, reporters, and restoration boundaries.}
\sourcefile{lib/report.abclr2}{sec:report}{Copy-based arrival/dispatch/completion reports and filtering.}
\end{longtable}


\section{Syntax at a Glance}
\label{sec:syntax}
This is a reading aid, not a complete grammar. Names such as \code{target},
\code{pattern}, and \code{body} are placeholders; ellipses stand for omitted
parts. The linked sections explain evaluation, matching, prerequisites,
and the limits of the supplied implementation. Lisp notation and control
forms are introduced in Section~\ref{sec:lisp-reading}.

\paragraph{Objects, state, and matching.}
\begin{longtable}{L{0.39\textwidth}L{0.39\textwidth}L{0.12\textwidth}}
\textbf{Construct} & \textbf{Meaning} & \textbf{Section}\\\hline
\endhead
\code{[object name (script ...)]} & Create an object and bind its name;
omit the name for an anonymous object. & \ref{sec:objects}\\
\code{(state [n := 0])} & Object-definition entry for persistent state. & \ref{sec:state}\\
\code{(temporary [x := 0])} & Handler-local binding, not persistent state. & \ref{sec:objects}\\
\code{[x := expression]} & Change a binding. & \ref{sec:state}\\
\code{(=> pattern body)} & A script handling a matching message. & \ref{sec:objects}\\
\code{(=> pattern @ R from S body)} & Additionally bind reply destination
and sender; either header field can be omitted. & \ref{sec:objects}\\
\code{[head . tail]} & Pattern binding the head and entire tail of a pair;
dot notation is not an extra list element. & \ref{sec:lisp-reading}\\
\code{(match value (is pattern body) ...)} & Match a value inside a script;
\code{otherwise} supplies a fallback clause. & \ref{sec:objects}\\
\code{(=> pattern where test body)} & Guarded script selection, requiring
compatible heavy-meta machinery; not an ordinary \code{match} constraint. & \ref{sec:guards}\\
\end{longtable}

\paragraph{Sending, replying, and continuing.}
\begin{longtable}{L{0.39\textwidth}L{0.39\textwidth}L{0.12\textwidth}}
\textbf{Construct} & \textbf{Meaning} & \textbf{Section}\\\hline
\endhead
\code{[target <= message]} & Past-type send; do not wait for a reply. & \ref{sec:messages}\\
\code{[target <== message]} & Now-type send; wait for a reply. & \ref{sec:messages}\\
\code{[target <= message @ R]} & Past-type send with a specified reply destination. & \ref{sec:messages}\\
\code{[target <= message \$ f]} & Future-type send; route replies through
the loaded future library. & \ref{sec:futures}\\
\code{!expression} & Send a reply for the current message; not a return
from the rest of the handler. & \ref{sec:objects}\\
\code{[cont pattern body]} & Create a short, non-reifying continuation object. & \ref{sec:messages}\\
\code{[target <= message :express]} & Express send; requires both sender
and receiver machinery, not just this extra parameter. & \ref{sec:message-extensions}\\
\code{(atomic body ...)} & Receiver-local protection in the express
protocol; not a multi-object transaction. & \ref{sec:message-extensions}\\
\end{longtable}

\noindent\begin{minipage}{\linewidth}
\paragraph{Groups, identities, and reflection.}
\begin{longtable}{L{0.39\textwidth}L{0.39\textwidth}L{0.12\textwidth}}
\textbf{Construct} & \textbf{Meaning} & \textbf{Section}\\\hline
\endhead
\code{[group name ...]} & Create a group manager with optional custom
kernel objects. & \ref{sec:reflection}\\
\code{(group manager)} & Object-definition parameter selecting its
creation group. & \ref{sec:reflection}\\
\code{(meta-gen generator)} & Object/group definition parameter selecting
a metaobject generator. & \ref{sec:reflection}\\
\code{(evaluator evaluator-object)} & Group-definition parameter selecting
its evaluator. & \ref{sec:reflection}\\
\code{[meta object]} & Obtain its metaobject. & \ref{sec:reflection}\\
\code{[den metaobject]} & Obtain its lower-level denotation, not a copy. & \ref{sec:reflection}\\
\code{[group-of object]} & Query its group manager; custom generators must
provide the corresponding identity metadata. & \ref{sec:original-inherit}\\
\code{Me} / \code{group} & Current object denoted by \code{Id} / group
manager in \code{Gid}, not necessarily base-level application objects. & \ref{sec:reflection}\\
\code{(reify-args tag forms ...)} & Evaluate the tag, then send compiled
argument expressions and context to the current evaluator. & \ref{sec:reflection}\\
\code{(reify-args-to-meta tag forms ...)} & Send them to the current
object's metaobject instead. & \ref{sec:reflection}\\
\end{longtable}
\end{minipage}

\code{reify-args} and \code{reify-args-to-meta} are unavailable in a
non-reifying object's script. Extra send/object/assignment parameters
are interpreted by the selected machinery, not universally by their
spelling. Section~\ref{sec:reflection} explains the evaluation context;
Section~\ref{sec:files} separately collects loading and compilation commands.


\section{Troubleshooting Program Protocols}
\label{sec:failures}

\begin{description}
  \item[No return messages.]
  A now-type call has no reply. Check that a script matched and every
  result path either uses \code{!} or forwards the original reply
  destination. Also look earlier in the transcript for a Lisp error.
  \item[An object is not bound.]
  Its definition may have failed while loading the file. Read the first
  error rather than treating later missing bindings as independent failures.
  \item[A future never finishes.]
  The producer may not send the terminal marker, or a reader may have
  consumed it already. Futures do not infer completion.
  \item[A macro expands into an unsupported special form.]
  ABCL/R2 does not support every Common Lisp special form. Inspect the
  macro expansion and compare it with the supported forms following
  Section~\ref{sec:lisp-reading}.
  \item[Script lambdas.]
  Lambdas embedded in scripts cannot capture ABCL lexical variables,
  and their bodies cannot contain ABCL object creation or message sends
  in this implementation~\cite{guide}. Prefer a script object for
  an ABCL callback.
\end{description}

\clearpage
\section{Concept and Protocol Index}
\label{sec:concept-index}

Use this alphabetical index to return to an explanation, not as a list of
every occurrence. Entries give the section and its starting page. For
notation use Appendix~\ref{sec:syntax}; for original filenames use
Appendix~\ref{sec:source-map}.

\begingroup
\renewcommand{\arraystretch}{1.4}
\newcommand{\indexroute}[1]{\ref{#1} (p.~\pageref{#1})}
\begin{longtable}{@{}p{0.52\textwidth}p{0.43\textwidth}@{}}
\textbf{Concept or protocol} & \textbf{Explanation}\\\hline
\endhead
Actors and asynchronous messages & \indexroute{sec:messages}\\
Assignment & \indexroute{sec:state}\\
Atomic regions and their limits & \indexroute{sec:message-extensions}; \indexroute{sec:original-express}\\
Bindings, lexical scope, and Lisp notation & \indexroute{sec:lisp-reading}\\
Breakpoints and environment inspection & \indexroute{sec:break}\\
Built-in object protocols & \indexroute{sec:builtin-reading}\\
Callbacks and script lambdas & \indexroute{sec:objects}; \indexroute{sec:failures}\\
Compilation and cached translations & \indexroute{sec:loading}\\
Completion: application versus evaluation & \indexroute{sec:originals}\\
Continuations, \code{cont}, and explicit reply targets & \indexroute{sec:messages}\\
Copy: identity, state, and shared references & \indexroute{sec:copy}\\
Denotation, \code{den}, and \code{meta} & \indexroute{sec:reflection}\\
Dispatch and script selection & \indexroute{sec:objects}; \indexroute{sec:guards}\\
Environment frames and inspection & \indexroute{sec:reflection}; \indexroute{sec:break}\\
Evaluation context: Env, Id, Gid, Eval, C & \indexroute{sec:reflection}\\
Express messages and preemption & \indexroute{sec:message-extensions}; \indexroute{sec:original-express}\\
Fault tolerance: rejection and retry & \indexroute{sec:retry}\\
Fibonacci: recursive objects and reply joins & \indexroute{sec:pcontrol-fib}\\
Futures, \code{make-future}, and \code{next-value} & \indexroute{sec:futures}; \indexroute{sec:original-future}\\
Groups, \code{group-of}, and object creation & \indexroute{sec:reflection}; \indexroute{sec:original-inherit}\\
Guards, \code{where}, and heavy-meta & \indexroute{sec:guards}\\
Host events and the pending-request slot & \indexroute{sec:host-events}\\
Identity and metaobject towers & \indexroute{sec:meta-meta-reading}\\
Inheritance and forwarding reply destinations & \indexroute{sec:original-inherit}\\
Interrupting evaluation & \indexroute{sec:interrupt}\\
Job-service protocol: run, square, done, report & \indexroute{sec:project}\\
Loading: \code{inc}, \code{comp}, and \code{inc-lib} & \indexroute{sec:loading}\\
Library search path: \code{*library-directories*} & \indexroute{sec:library-path}\\
Macros and supported Lisp forms & \indexroute{sec:first-session}; \indexroute{sec:failures}\\
Matching, dotted patterns, and fallback & \indexroute{sec:objects}\\
MCP and the meta-circular interpreter & \indexroute{sec:mcp-reading}\\
Me and the currently executing object & \indexroute{sec:reflection}\\
Messages: past-, now-, and future-type & \indexroute{sec:messages}; \indexroute{sec:futures}\\
Modules: \code{require} and \code{provide} & \indexroute{sec:loading}\\
Monitor measurements and their limits & \indexroute{sec:pcontrol-monitor}\\
Multicast and positional reply collection & \indexroute{sec:multicast}\\
Non-reifying execution & \indexroute{sec:reflection}; \indexroute{sec:tw-variants}\\
Objects, factories, and scripts & \indexroute{sec:objects}; \indexroute{sec:account}\\
Priority queues and scheduling policy & \indexroute{sec:pcontrol-policy}\\
Quicksort as a message-driven state machine & \indexroute{sec:pcontrol-qsort}\\
Reflection: evaluator and metaobject routes & \indexroute{sec:reflection}; \indexroute{sec:original-if}\\
Reification: \code{reify-args}, \code{reify-args-to-meta} & \indexroute{sec:reflection}; \indexroute{sec:interrupt}\\
Replies, \code{!}, sender, and destination & \indexroute{sec:objects}; \indexroute{sec:originals}\\
Report: receiving-side observation and restoration & \indexroute{sec:report}\\
Rollback, anti-messages, and replay & \indexroute{sec:tw-recovery}; \indexroute{sec:tw-cascade}\\
Script installation and dynamic updates & \indexroute{sec:user-script-reading}\\
State versus temporary bindings & \indexroute{sec:state}; \indexroute{sec:objects}\\
Streams and terminal-message protocols & \indexroute{sec:original-future}; \indexroute{sol:stream}\\
Time Warp: virtual time and simulation & \indexroute{sec:timewarp}\\
Transactions: limits of an account transfer & \indexroute{sol:transfer}\\
Watch: outgoing-send observation and restoration & \indexroute{sec:watch}\\
\end{longtable}
\endgroup

